\begin{helper}
Let $\mathcal F$ be a finite multi-hypergraph, let $f$ be an edge, and let
$v\in f$.  In the fresh-edge hypergraph
$\operatorname{FreshEdge}_k(\mathcal F,v)$ of
\noderef{KUniformFreshEdgeHypergraph}, write $\operatorname{inl}(f)$ for the
old-edge copy of $f$ and write $\operatorname{inr}(*)$ for the unique fresh
edge label.  Then the number of line-graph neighbors of $\operatorname{inl}(f)$
in $L(\operatorname{FreshEdge}_k(\mathcal F,v))$ is the number of line-graph
neighbors of $f$ in $L(\mathcal F)$ plus one.
\end{helper}
\begin{proof}
By \noderef{LineGraphOfHypergraph}, an edge is adjacent to $f$ in the line
graph exactly when it is different from $f$ and its hyperedge has nonempty
intersection with the hyperedge of $f$.

Consider first an old edge $g$ of $\mathcal F$.  In
\noderef{KUniformFreshEdgeHypergraph}, the copy of $g$ has edge set
$g$ with every vertex sent by the old-vertex inclusion.  Therefore the copy
of $g$ meets the copy of $f$ exactly when $g$ met $f$ before.  Also the old
edge labels remain distinct under the old-edge inclusion.  Hence the old
neighbors of $f$ are in bijection with the old-edge copies that neighbor the
copy of $f$ in the fresh-edge hypergraph.

There is exactly one non-old edge label in the fresh-edge construction.  Its
edge contains the old vertex $v$ and the fresh vertices.  Since $v\in f$, this
new edge has nonempty intersection with the copy of $f$, and it is plainly not
the old label $f$.  Thus it contributes one additional line-graph neighbor.
No old-edge copy is equal to this new edge label, so the old-neighbor image
and the singleton consisting of the new edge label are disjoint.  Taking
finite cardinalities gives the asserted increase by one.
\end{proof}
